\documentclass[../candidatura.tex]{subfiles}
\graphicspath{{\subfix{../imagenes/}}}
\begin{document}

\section{Popusintesíntesis sonora}

\avisoesbozo

La componente popular se propone como un conjunto de estrategias de diseño y producción, aplicada a la síntesis de sintetizadores sonoros.

Este prefijo popular es un homenaje a la serie volca de Tatsuya Takahashi en KORG y a los Populibros Peruanos de Manuel Scorza, descritos en el marco teórico (sección \ref{subsec:marco-teorico}).

\subsection{Legibilidad}

Que un instrumento sea legible significa que su interfaz permite a una persona usuaria comprender cómo usarlo, cómo se relacionan sus parámetros entre sí.

También permite a una entusiasta o a una artesana comprender cómo se construye, cómo se repara, cómo se modifica. 

\subsection{Reparabilidad}

Se propone que los instrumentos populares están hechos a partir de componentes asequibles, sin escasez, y que su diseño se hace cargo del ciclo de vida del instrumento, incluyendo estrategias para repararlo, clonarlo, modificarlo.


\newpage

\end{document}
